% Source: https://artofproblemsolving.com/community/c3h1851993_inverse_functions
% Stable identity: aops-topic-1851993-q01
\textbf{(a)} Let $g(x)=e^x$, and let $f(x)=\ln|x|$ for $x\ne0$ and $f(0)=0$. Then $f(g(x))=\ln e^x=x$ for every real $x$, but
\[
g(f(-1))=e^{\ln1}=e^0=1\ne-1 .
\]
So $g(f(x))=x$ fails, for example at $x=-1$ (in fact for every $x\le0$). The reason is that $g$ only takes positive values, and $f$ can do anything on the numbers that $g$ never reaches.

\textbf{(b)} If $g(s)=g(t)$, then $s=f(g(s))=f(g(t))=t$, so $g$ is one-to-one. Every real $y$ equals $f(g(y))$, so $y$ is a value of $f$; hence $f$ is onto.

\textbf{(c)} Let $x$ be real. Since $g$ is onto, $x=g(t)$ for some $t$. Then
\[
g(f(x))=g\bigl(f(g(t))\bigr)=g(t)=x .\qquad\blacksquare
\]
So the answer to ``can $f(g(x))=x$ hold while $g(f(x))\ne x$?'' is \textbf{yes}, and this happens exactly when $g$ is not onto.
